% ----------------------------------------------------------------------
%  3.3. Теория фаз в сплавах: 3.3.1 классификация, 3.3.2 твёрдые растворы
% ----------------------------------------------------------------------
\begin{orientir}
Что получится, если сплавить два металла? Атомы второго компонента
могут занять места атомов первого, втиснуться между ними или вместе с
ними построить совсем новую решётку. В этом разделе мы разберём
\emph{твёрдые растворы} (3.3.2)~— замещения, внедрения и вычитания,
условия неограниченной растворимости и упорядочение,~— и
\emph{промежуточные фазы} (3.3.3), прежде всего электронные соединения
и фазы внедрения. Здесь пригодятся правило отношения радиусов
(п.~\ref{sssec:ionic}) и зонная теория (п.~\ref{ssec:bands}).
\end{orientir}

\subsection{Классификация фаз в сплавах}\label{ssec:phase-class}

\begin{opredelenie}
\term{Фаза}~— химически и физически однородная часть неоднородной
системы, отделённая от других частей поверхностью раздела.
\term{Сплав} (по Коттреллу)~— макроскопически однородное вещество с
металлическими свойствами, состоящее из двух или более химических
элементов.
\end{opredelenie}

Для простоты мы будем рассматривать в основном \emph{двойные} (бинарные)
сплавы. При данных давлении, температуре и составе сплав может быть
однофазным (гомогенным) или многофазным (гетерогенным).

Общепринятая классификация фаз опирается на их положение на диаграмме
состояния и на кристаллическую структуру. С этой точки зрения твёрдые
фазы делят на два типа:
\begin{itemize}
\item \term{первичные твёрдые растворы}~— примыкают на диаграмме к
чистым компонентам и имеют решётку одного из них;
\item \term{промежуточные фазы}~— все остальные; они не прилегают к
чистым компонентам и обычно имеют собственную решётку.
\end{itemize}
Реже фазы классифицируют по степени упорядоченности атомов или по
электронному строению. Последняя классификация была бы самой
глубокой, но точных моделей электронного строения многоэлектронных
сплавов нет, и концентрацию валентных электронов часто нельзя
определить однозначно. Удачные исключения~— электронные соединения
(п.~\ref{ssec:intermediate}).

\subsection{Твёрдые растворы}\label{ssec:solid-sol}

\begin{opredelenie}
\term{Твёрдый раствор}~— твёрдая фаза, в которой соотношение
компонентов может меняться без нарушения однородности и с сохранением
одной и той же кристаллической решётки.
\end{opredelenie}

\paragraph{Три типа твёрдых растворов}
По тому, где располагаются атомы второго компонента, различают растворы
замещения, внедрения и вычитания (рис.~\ref{fig:ss-types}).

\begin{figure}[ht]
\centering
\begin{tikzpicture}[scale=0.5]
  % а — замещение
  \begin{scope}
    \foreach \x in {0,...,5}{\foreach \y in {0,...,5}{\node[atom=3.6mm] at (\x,\y) {};}}
    \foreach \p in {(1,4),(3,3),(4,1),(2,1),(5,4)}{\node[atomB=3.6mm] at \p {};}
    \node[font=\small] at (2.5,-1.1) {\textit{а}};
  \end{scope}
  % б — внедрение
  \begin{scope}[shift={(7.5,0)}]
    \foreach \x in {0,...,5}{\foreach \y in {0,...,5}{\node[atom=3.6mm] at (\x,\y) {};}}
    \foreach \p in {(1.5,3.5),(3.5,1.5),(2.5,4.5),(4.5,2.5)}{\node[atomI=2mm] at \p {};}
    \node[font=\small] at (2.5,-1.1) {\textit{б}};
  \end{scope}
  % в — вычитание (соединение AB с вакансиями в подрешётке A)
  \begin{scope}[shift={(15,0)}]
    \foreach \x in {0,...,5}{\foreach \y in {0,...,5}{
      \pgfmathtruncatemacro{\prt}{mod(\x+\y,2)}
      \ifnum\prt=0 \node[atom=3.6mm] at (\x,\y) {};\else\node[atomB=3.0mm] at (\x,\y) {};\fi}}
    \foreach \p in {(1,2),(4,3),(2,5),(3,0)}{\node[vac=3.6mm] at \p {}; \node[font=\scriptsize, fmRed] at \p {$\times$};}
    \node[font=\small] at (2.5,-1.1) {\textit{в}};
  \end{scope}
\end{tikzpicture}
\caption{Типы твёрдых растворов: \textit{а}~— замещения (оранжевые
атомы второго компонента занимают узлы решётки растворителя);
\textit{б}~— внедрения (малые атомы в междоузлиях); \textit{в}~—
вычитания (на основе соединения; $\times$~— незанятые узлы одной из
подрешёток)}
\label{fig:ss-types}
\end{figure}

\begin{itemize}
\item \term{Растворы замещения}: атомы растворённого элемента
\emph{замещают} атомы растворителя в узлах решётки, распределяясь
хаотически. Так растворяются друг в друге и химические соединения:
NaCl--KCl, TiC--VC, TiC--TiN и даже карбонитрид (Ti,\,V)(C,\,N). Здесь
атомы металлов и неметаллов замещают друг друга \emph{каждый в своей
подрешётке}.
\item \term{Растворы внедрения}: малые атомы (H, C, N, O, а также B)
располагаются в \emph{междоузлиях}~— пустотах решётки растворителя.
\item \term{Растворы вычитания} (с дефектной решёткой) образуются только
на основе химических соединений: часть позиций одного из компонентов
остаётся \emph{незанятой}.
\end{itemize}

\paragraph{Растворы вычитания и переменная валентность}
Пример~— \term{вюстит} FeO со структурой NaCl: все позиции кислорода
заняты, а часть позиций железа пуста. В нитридах циркония и гафния с
избытком азота пусты часть позиций металла. Как сохраняется
электронейтральность? Только за счёт \emph{разной валентности} ионов
одного элемента. Кислород в оксидах всегда двухвалентен; меняется
валентность ионов переходного металла (у них близки уровни $ns$ и
$(n-1)d$). В вюстите наряду с Fe$^{2+}$ появляются ионы Fe$^{3+}$.

Ярко это видно у оксидов титана со структурой NaCl
(табл.~\ref{tab:TiO}). Вакансии здесь бывают в обеих подрешётках.

\begin{table}[ht]
\centering
\caption{Вакансии в оксидах титана со структурой NaCl}
\label{tab:TiO}
\small
\begin{tabular}{@{}lccc@{}}
\toprule
 & TiO$_{0,6}$ & TiO & Ti$_{0,8}$O \\
\midrule
кислород, ат.\,\% & $\frac{60}{160}=38$ & 50 & $\frac{100}{180}=56$ \\
вакансии в подрешётке Ti & 0 & 15\,\% & 20\,\% \\
вакансии в подрешётке O & 40\,\% & 15\,\% & 0 \\
доля вакантных узлов (всего) & 20\,\% & 15\,\% & 10\,\% \\
\bottomrule
\end{tabular}
\end{table}

При недостатке кислорода в оксиде появляются ионы титана пониженной
валентности, при избытке~— четырёхвалентные.

\paragraph{Искажения решётки}
Атомы растворителя и растворённого элемента различаются размерами, и
решётка твёрдого раствора искажается (рис.~\ref{fig:ss-distort}). Ещё
сильнее искажения при внедрении атомов в междоузлия или при образовании
вакансий. Искажения подтверждаются тремя группами фактов:
\begin{enumerate}
\item на рентгенограммах интенсивность линий (особенно с большими
индексами) ниже, чем у чистого растворителя, а фон выше;
\item электросопротивление растёт (рис.~\ref{fig:ss-props},\,\textit{а});
\item твёрдость растёт, а пластичность падает (рис.~\ref{fig:ss-props},\,\textit{б}).
\end{enumerate}

\begin{figure}[ht]
\centering
\begin{tikzpicture}[scale=0.62]
  \foreach \sh/\s/\col/\sz/\lab in {0/0.28/atomB/5.4mm/{\textit{а}~— атом крупнее},
                                      8.2/-0.22/atomB/2.6mm/{\textit{б}~— атом мельче}}{
    \begin{scope}[shift={(\sh,0)}]
      \foreach \i in {-3,...,3}{\foreach \j in {-3,...,3}{
        \pgfmathsetmacro{\rr}{\i*\i+\j*\j}
        \ifdim\rr pt>0pt
          \pgfmathsetmacro{\xx}{\i+\s*\i/\rr}
          \pgfmathsetmacro{\yy}{\j+\s*\j/\rr}
          \node[atom=3.8mm] at (\xx,\yy) {};
        \fi}}
      \node[atom=\sz, fill=\col] at (0,0) {};
      \foreach \k in {-3,...,3}{\draw[fmGray!40, line width=0.3pt, dashed] (\k,-3.6) -- (\k,3.6) (-3.6,\k) -- (3.6,\k);}
      \node[font=\small] at (0,-4.4) {\lab};
    \end{scope}}
\end{tikzpicture}
\caption{Искажение решётки вокруг растворённого атома (расчёт по
упругому полю смещений $u\propto1/r$; смещения увеличены для
наглядности). Пунктир~— положения узлов идеальной решётки:
соседи раздвигаются вокруг крупного атома и стягиваются к мелкому}
\label{fig:ss-distort}
\end{figure}

\begin{figure}[ht]
\centering
\begin{subfigure}[b]{0.48\textwidth}\centering
\begin{tikzpicture}
\begin{axis}[width=0.9\textwidth, height=5cm, xmin=0, xmax=100, ymin=0, ymax=55,
  xlabel={Ni, ат.\,\%}, ylabel={$\rho$, мкОм$\cdot$см}, axis lines=left]
  \addplot[fmBlue, thick, domain=0:100, samples=80] {1.7*(1-x/100)+6.9*x/100+175*(x/100)*(1-x/100)};
  \addplot[fmGray, dashed, domain=0:100] {1.7*(1-x/100)+6.9*x/100};
  \node[font=\scriptsize, fmGray, anchor=south east] at (axis cs:98,9) {аддитивная прямая};
\end{axis}
\end{tikzpicture}
\caption{}
\end{subfigure}\hfill
\begin{subfigure}[b]{0.48\textwidth}\centering
\begin{tikzpicture}
\begin{axis}[width=0.9\textwidth, height=5cm, xmin=0, xmax=100, ymin=0, ymax=1.3,
  xlabel={B, ат.\,\%}, ylabel={твёрдость $H_V$}, ytick=\empty, axis lines=left]
  \addplot[fmRed, thick, domain=0:100, samples=80] {0.4*(1-x/100)+0.6*x/100+1.8*(x/100)*(1-x/100)};
  \addplot[fmGray, dashed, domain=0:100] {0.4*(1-x/100)+0.6*x/100};
  \node[font=\scriptsize] at (axis cs:8,0.3) {A}; \node[font=\scriptsize] at (axis cs:93,0.5) {B};
\end{axis}
\end{tikzpicture}
\caption{}
\end{subfigure}
\caption{Изменение свойств при образовании непрерывного ряда твёрдых
растворов: \textit{а}~— удельное электросопротивление сплавов Cu--Ni
(оценка по правилу Нордгейма $\rho\propto c(1-c)$ с опорой на
сопротивление чистых металлов и максимум около 50~мкОм$\cdot$см);
\textit{б}~— твёрдость (схема): максимумы при промежуточных составах}
\label{fig:ss-props}
\end{figure}

Из-за искажений период решётки твёрдого раствора в разных местах
кристалла разный~— в зависимости от локального состава. Рентгеновский
анализ даёт лишь \emph{среднее} значение, которое и называют периодом
решётки твёрдого раствора. Как он меняется с составом, зависит от типа
раствора (рис.~\ref{fig:a-vs-c}):
\begin{itemize}
\item в растворах \emph{внедрения} период всегда \emph{растёт}
(аустенит~— раствор углерода в $\gamma$-Fe);
\item в растворах \emph{вычитания} период с ростом числа вакансий
\emph{уменьшается}. Стехиометрическому FeO отвечает 22,3~мас.\,\%
кислорода; реальный вюстит содержит больше кислорода, т.\,е. меньше
железа, и чем больше вакансий железа, тем меньше период;
\item в растворах \emph{замещения} период может и расти, и уменьшаться в
зависимости от размеров растворённых атомов.
\end{itemize}

\begin{figure}[ht]
\centering
\begin{subfigure}[b]{0.48\textwidth}\centering
\begin{tikzpicture}
\begin{axis}[width=0.9\textwidth, height=5cm, xmin=0, xmax=2, ymin=0.3570, ymax=0.3650,
  xlabel={C, мас.\,\%}, ylabel={$a$, нм}, axis lines=left,
  yticklabel style={/pgf/number format/fixed, /pgf/number format/precision=4},
  scaled y ticks=false]
  \addplot[fmBlue, thick, domain=0:2] {0.3578+0.0033*x};
\end{axis}
\end{tikzpicture}
\caption{}
\end{subfigure}\hfill
\begin{subfigure}[b]{0.48\textwidth}\centering
\begin{tikzpicture}
\begin{axis}[width=0.9\textwidth, height=5cm, xmin=23.0, xmax=24.6, ymin=0.4270, ymax=0.4315,
  xlabel={O, мас.\,\%}, ylabel={$a$, нм}, axis lines=left,
  yticklabel style={/pgf/number format/fixed, /pgf/number format/precision=4},
  xticklabel style={/pgf/number format/fixed},
  scaled y ticks=false]
  % Fe(1-y)O: мас. доля O = 16/(16+(1-y)·55,85); a = 0,4334 - 0,0478·y
  \addplot[fmRed, thick, domain=0.04:0.125, samples=30] ({1600/(16+(1-x)*55.85)}, {0.4334-0.0478*x});
\end{axis}
\end{tikzpicture}
\caption{}
\end{subfigure}
\caption{Изменение периода решётки с составом (по эмпирическим
зависимостям): \textit{а}~— раствор внедрения (аустенит): период растёт
с содержанием углерода; \textit{б}~— раствор вычитания (вюстит
Fe$_{1-y}$O): с ростом доли кислорода растёт число вакансий железа, и
период уменьшается}
\label{fig:a-vs-c}
\end{figure}

\paragraph{Закон Вегарда}
Если период решётки твёрдого раствора меняется по прямой между
периодами чистых компонентов $a_1$ и $a_2$, говорят, что выполняется
\term{закон Вегарда}:
\begin{keyeq}\label{eq:vegard}
a=a_1+(a_2-a_1)\,c ,
\end{keyeq}
где $c$~— атомная доля второго компонента. В металлических системах
этот закон строго не выполняется (рис.~\ref{fig:vegard}): у Cu--Ni
кривая слегка «провисает» (отрицательное отклонение), у Cu--Au
отклонение положительное, а у Au--Pt и Au--Pd закон выполняется почти
точно\footnote{В оригинале знак отклонения связывается с разницей
атомных диаметров компонентов (меньше или больше 9\,\%). Однако
приведённые там же примеры (Cu--Pd, Au--Pt, Au--Pd с разницей диаметров
около 8, 4 и 5\,\%) этому правилу не соответствуют, поэтому мы его не
приводим.}. В растворах с ограниченной растворимостью отклонения
бывают обоих знаков; встречаются даже случаи, когда при растворении
\emph{меньших} атомов период \emph{растёт} (никель с $r=0{,}1244$~нм в
вольфраме с $r=0{,}1408$~нм). У сплавов $\alpha$-Fe--Cr кривая имеет
точку перегиба, у Cr--Ti~— три. Причины отклонений~— различие упругих
свойств компонентов (модулей, сжимаемости) и перераспределение
электронов между атомами.

\begin{figure}[ht]
\centering
\begin{tikzpicture}
\begin{axis}[width=0.62\textwidth, height=5.4cm, xmin=0, xmax=1, ymin=0.350, ymax=0.412,
  xlabel={атомная доля второго компонента $c$}, ylabel={$a$, нм}, axis lines=left,
  yticklabel style={/pgf/number format/fixed, /pgf/number format/precision=3}, scaled y ticks=false]
  \addplot[fmGray, dashed, domain=0:1] {0.3615+(0.4078-0.3615)*x};
  \addplot[fmRed, thick, domain=0:1, samples=50] {0.3615+(0.4078-0.3615)*x+0.012*x*(1-x)};
  \addplot[fmGray, dashed, domain=0:1] {0.3615+(0.3524-0.3615)*x};
  \addplot[fmBlue, thick, domain=0:1, samples=50] {0.3615+(0.3524-0.3615)*x-0.002*x*(1-x)};
  \node[font=\small, fmRed, anchor=south east] at (axis cs:0.95,0.404) {Cu--Au};
  \node[font=\small, fmBlue, anchor=north east] at (axis cs:0.95,0.3565) {Cu--Ni};
  \node[font=\scriptsize, fmGray, anchor=north west] at (axis cs:0.42,0.381) {закон Вегарда};
\end{axis}
\end{tikzpicture}
\caption{Отклонения от закона Вегарда (схема; крайние точки~— периоды
чистых Cu, Au и Ni). Пунктир~— линейный закон \eqref{eq:vegard}}
\label{fig:vegard}
\end{figure}

\paragraph{Как определить тип твёрдого раствора}
Нужно точно измерить период решётки $a$ и плотность $\rho$ и сосчитать,
сколько атомов приходится на элементарную ячейку:
\begin{equation}\label{eq:n-cell}
n=\frac{\rho\,V_\text{яч}}{m_u\,A_\text{ср}},
\end{equation}
где $V_\text{яч}$~— объём ячейки, $m_u=1{,}66\cdot10^{-24}$~г~— атомная
единица массы, $A_\text{ср}$~— средняя атомная масса сплава:
\[
A_\text{ср}=\frac{c_1A_1+c_2A_2+\ldots}{100}\ \ (c_i\text{ в ат.\,\%}),\qquad
A_\text{ср}=\frac{100}{c_1/A_1+c_2/A_2+\ldots}\ \ (c_i\text{ в мас.\,\%}).
\]
Пусть $n_0$~— число атомов в ячейке металла-основы. Тогда
\begin{center}
$n=n_0$~— раствор замещения;\quad $n>n_0$~— внедрения;\quad $n<n_0$~— вычитания.
\end{center}

\begin{primer}[Вюстит]
Состав 47,7~ат.\,\% Fe и 52,3~ат.\,\% O, плотность $\rho=5{,}613$~г/см$^3$,
решётка типа NaCl с $a=0{,}4282$~нм ($n_0=8$ ионов на ячейку).
\begin{shagi}
\item $A_\text{ср}=0{,}477\cdot55{,}85+0{,}523\cdot16{,}00=35{,}0$.
\item $V_\text{яч}=a^3=(0{,}4282\cdot10^{-7}~\text{см})^3=7{,}85\cdot10^{-23}$~см$^3$.
\item $n=\dfrac{5{,}613\cdot7{,}85\cdot10^{-23}}{1{,}66\cdot10^{-24}\cdot35{,}0}\approx7{,}6<8$.
\end{shagi}
Это раствор вычитания\isprmark. Дефицит около $0{,}36$ иона на ячейку составляет
4,5\,\% от восьми~— и совпадает с разницей содержаний кислорода и железа
($52{,}3-47{,}7=4{,}6$~ат.\,\%).
\end{primer}
\isprtext{В оригинале (с.~407) указано
$A_\text{ср}=36{,}0$~— это опечатка: при таком значении получилось бы
$n\approx7{,}4$. Приведённое там же $n=7{,}64$ соответствует
$A_\text{ср}=35{,}0$.}

\paragraph{Неограниченная растворимость}
Металлы могут растворяться друг в друге неограниченно только по типу
замещения. Для этого необходимы (но не достаточны!) три условия.
\begin{shagi}
\item \emph{Одинаковые кристаллические решётки} компонентов. При
полиморфизме неограниченная растворимость возможна только между
изоморфными модификациями.
\item \emph{Близкие атомные диаметры} (правило Юм-Розери): при одном и
том же координационном числе они должны различаться не более чем на
14--15\,\%.
\item \emph{Близкие электрохимические свойства}: близкие
электроотрицательности и электронное строение, т.\,е. одна или соседние
группы периодической системы. Иначе возникает тенденция к образованию
промежуточных фаз.
\end{shagi}

\emph{Щелочные металлы} наглядно показывают роль размерного фактора:
первое и третье условия выполнены автоматически (все~— ОЦК, все~— в
первой группе). Размеры же различаются сильно (табл.~\ref{tab:alkali}).
Неограниченная растворимость наблюдается только у K, Rb, Cs~— соседей с
разницей размеров 5--8\,\%. Натрий с ними твёрдых растворов не
образует, а литий из-за огромной разницы размеров не смешивается с
остальными щелочными металлами даже в жидком состоянии.

\begin{table}[ht]
\centering
\caption{Периоды решётки щелочных металлов и рост атомного диаметра}
\label{tab:alkali}
\small
\begin{tabular}{@{}lccccccccc@{}}
\toprule
 & Li & & Na & & K & & Rb & & Cs \\
\midrule
$a$, нм & 0,350 & & 0,428 & & 0,533 & & 0,562 & & 0,605 \\
рост диаметра, \% & & 22 & & 25 & & 5,5 & & 7,6 & \\
\bottomrule
\end{tabular}
\end{table}

\emph{Металлы группы меди}. Диаметр атома меди 0,255~нм, серебра~— на
13\,\% больше (0,2884~нм), а у более тяжёлого золота из-за лантаноидного
сжатия~— практически такой же, как у серебра (0,2879~нм). Au--Ag и Au--Cu
образуют непрерывные ряды твёрдых растворов, а Cu--Ag~— только
ограниченные: условия выполнены «на грани», и их оказывается
недостаточно.

Недостаточность трёх условий видна и на других примерах. Медь
неограниченно растворяется с никелем, а с железом, атом которого почти
такой же, как у никеля, взаимная растворимость в твёрдом состоянии лишь
единицы процентов\ispr{В оригинале (с.~409) сказано, что медь и железо
не смешиваются даже в жидком состоянии, и так же ведёт себя система
Cu--Ca. По современным диаграммам состояния жидкие Cu и Fe смешиваются
неограниченно (расслоение возможно лишь в сильно переохлаждённом
расплаве), а Cu и Ca образуют интерметаллиды (CaCu$_5$ и др.).}. Золото,
как и медь, неограниченно растворяется с никелем, а с кобальтом~—
ограниченно. Серебро же, по размеру почти равное золоту, не смешивается
ни с кобальтом, ни с никелем ни в твёрдом, ни даже в жидком состоянии.

\paragraph{Ограниченная растворимость}
Ограниченные растворы замещения образуют и металлы с разными решётками,
но размерный и электрохимический факторы действуют и здесь: чем они
благоприятнее, тем шире растворимость. Даркен и Гурри предложили
наглядный способ (рис.~\ref{fig:darken}): элементы наносят на карту
«электроотрицательность~— атомный диаметр» (при к.\,ч.~12) и вокруг
растворителя строят эллипс с полуосями $\pm0{,}4$ по
электроотрицательности и $\pm15$\,\% по диаметру. Элементы внутри
эллипса, как правило, дают широкие твёрдые растворы (более 5~ат.\,\%),
вне его~— очень узкие.

\begin{figure}[ht]
\centering
\begin{tikzpicture}[x=2.6cm, y=2.6cm]
  \draw[-{Stealth[length=2mm]}] (0,0) -- (2.5,0) node[below left, font=\small]{атомный диаметр};
  \draw[-{Stealth[length=2mm]}] (0,0) -- (0,1.75) node[above, font=\small]{электроотрицательность $\chi$};
  \fill[fmGreen!15] (1.25,0.85) ellipse (0.45 and 0.4);
  \draw[fmGreen!60!black, thick] (1.25,0.85) ellipse (0.45 and 0.4);
  \fill[fmBlue] (1.25,0.85) circle (2.2pt) node[below right, font=\small]{растворитель};
  \draw[dimline] (1.25,1.32) -- node[above, font=\scriptsize]{$\pm15$\,\%} (1.7,1.32);
  \draw[dimline] (1.78,0.85) -- node[right, font=\scriptsize]{$\pm0{,}4$} (1.78,1.25);
  \foreach \p in {(1.0,0.7),(1.45,1.05),(1.15,1.05),(1.4,0.62)}{\fill[fmGreen!50!black] \p circle (1.6pt);}
  \foreach \p in {(0.5,0.9),(2.0,0.5),(1.3,1.5),(2.1,1.3),(0.65,0.35)}{\draw[fmRed, line width=0.7pt] \p ++(-0.03,-0.03) -- ++(0.06,0.06) ++(-0.06,0) -- ++(0.06,-0.06);}
  \node[font=\scriptsize, fmGreen!50!black, align=center] at (1.25,0.28) {широкая\\растворимость};
  \node[font=\scriptsize, fmRed] at (2.15,0.38) {узкая};
\end{tikzpicture}
\caption{Принцип диаграммы Даркена--Гурри (схема): элементы внутри эллипса
(точки) образуют с растворителем широкие твёрдые растворы, вне его
(крестики)~— очень узкие}
\label{fig:darken}
\end{figure}

Ещё одна закономерность: металл с \emph{большей} валентностью обычно
растворяет металл с меньшей валентностью хуже, чем растворяется в нём
сам. Серебро растворяет до 30\,\% Mg, а магний~— лишь 3,9\,\% Ag; медь
растворяет до 6,5\,\% Mg, а магний~— не более 0,01\,\% Cu.

\paragraph{Электронная концентрация}
На растворимость влияет и \term{электронная концентрация}~— число
валентных электронов на атом сплава, $e/a$. Возьмём металлы, следующие
за медью в четвёртом периоде: Zn, Ga, Ge, As (валентность $z=2\ldots5$).
Чем больше валентность, тем меньше их растворимость в меди
(табл.~\ref{tab:Cu-sol}). Но если пересчитать предельные концентрации в
электронную концентрацию,
\[
e/a=1\cdot(1-c)+z\,c=1+(z-1)\,c ,
\]
то для Zn, Ga, Ge получается почти одно и то же число~— около 1,36--1,39
(рис.~\ref{fig:Cu-sol}). Иначе говоря, растворение прекращается, когда
электронная концентрация достигает предельного значения около 1,36.
Это число мы уже встречали: при $e/a=1{,}36$ сфера Ферми касается
границы зоны Бриллюэна ГЦК-решётки (п.~\ref{ssec:bands}). Почему это
важно, объясним в п.~\ref{ssec:intermediate}.

\begin{table}[ht]
\centering
\caption{Растворимость элементов четвёртого периода в меди}
\label{tab:Cu-sol}
\small
\begin{tabular}{@{}lcccc@{}}
\toprule
Система & Валентность $z$ & $\Delta d/d$, \% & $c_\text{max}$, ат.\,\% & $(e/a)_\text{max}$ \\
\midrule
Cu--Zn & 2 & $+7{,}2$ & 38,4 & 1,38 \\
Cu--Ga & 3 & $+6{,}6$ & 19,3 & 1,39 \\
Cu--Ge & 4 & $+8{,}5$ & 12,0 & 1,36 \\
Cu--As & 5 & $+15{,}0$ & 6,8 & 1,27 \\
\bottomrule
\end{tabular}
\end{table}

\begin{figure}[ht]
\centering
\begin{tikzpicture}
\begin{axis}[width=0.6\textwidth, height=5.2cm, xmin=1.6, xmax=5.4, ymin=0, ymax=45,
  xlabel={валентность растворённого элемента $z$}, ylabel={$c_\text{max}$, ат.\,\%},
  xtick={2,3,4,5}, axis lines=left]
  \addplot[fmGray, dashed, domain=1.7:5.3, samples=60] {36/(x-1)};
  \addplot[only marks, mark=*, mark size=2.2pt, fmBlue] coordinates {(2,38.4) (3,19.3) (4,12.0) (5,6.8)};
  \node[font=\scriptsize, fmBlue, anchor=west] at (axis cs:2.08,38.4) {Zn};
  \node[font=\scriptsize, fmBlue, anchor=south west] at (axis cs:3.05,19.3) {Ga};
  \node[font=\scriptsize, fmBlue, anchor=south west] at (axis cs:4.05,12.0) {Ge};
  \node[font=\scriptsize, fmBlue, anchor=north east] at (axis cs:4.95,6.6) {As};
  \node[font=\scriptsize, fmGray, align=left, anchor=west] at (axis cs:3.4,30) {$e/a=1{,}36$:\\$c=\frac{0{,}36}{z-1}$};
\end{axis}
\end{tikzpicture}
\caption{Предельная растворимость в меди и правило электронной
концентрации: пунктир~— составы, при которых $e/a=1{,}36$. Мышьяк
растворяется меньше предсказанного: у него неблагоприятен ещё и
размерный фактор (15\,\%)}
\label{fig:Cu-sol}
\end{figure}

\paragraph{Упорядоченные твёрдые растворы (сверхструктуры)}
Искажения решётки заметно уменьшаются, если атомы компонентов
располагаются в ней \emph{упорядоченно}. Впервые такой эффект обнаружили
Н.\,С.~Курнаков с сотрудниками на сплавах золота с медью. Отжиг
закалённых (с температур выше 500\,°C) сплавов составов Cu$_3$Au и CuAu
уменьшал их электросопротивление в 2--3 раза, а при длительном
ступенчатом отжиге с медленным (до 10 суток) понижением температуры~—
ещё сильнее (рис.~\ref{fig:CuAu-rho}). У закалённого
(неупорядоченного) раствора сопротивление максимально при
эквиатомном составе; у отожжённого сплава того же состава оно падает
почти до уровня чистых металлов.

\begin{figure}[ht]
\centering
\begin{tikzpicture}
\begin{axis}[width=0.62\textwidth, height=5.2cm, xmin=0, xmax=100, ymin=0, ymax=17,
  xlabel={Au, ат.\,\%}, ylabel={$\rho$, мкОм$\cdot$см}, axis lines=left, xtick={0,25,50,75,100}]
  \addplot[fmRed, thick, domain=0:100, samples=80] {1.7*(1-x/100)+2.2*x/100+56*(x/100)*(1-x/100)};
  \addplot[fmBlue, thick, smooth] coordinates {(0,1.7) (8,5.5) (15,7.8) (20,6.5) (25,4.2) (30,6.8) (38,9.5) (45,7.0) (50,4.5) (55,7.5) (62,11.5) (75,10.5) (88,6.5) (100,2.2)};
  \node[font=\small, fmRed, anchor=south west] at (axis cs:66,14.6) {закалка (беспорядок)};
  \node[font=\small, fmBlue, anchor=west] at (axis cs:78,3.2) {отжиг (порядок)};
  \draw[fmGray, dotted] (axis cs:25,0) -- (axis cs:25,4.2);
  \draw[fmGray, dotted] (axis cs:50,0) -- (axis cs:50,4.5);
  \node[font=\scriptsize, anchor=south] at (axis cs:25,0.2) {Cu$_3$Au};
  \node[font=\scriptsize, anchor=south] at (axis cs:50,0.2) {CuAu};
\end{axis}
\end{tikzpicture}
\caption{Электросопротивление сплавов Cu--Au в закалённом и отожжённом
состояниях (схема по опытам Курнакова): упорядочение при составах
Cu$_3$Au и CuAu резко снижает сопротивление}
\label{fig:CuAu-rho}
\end{figure}

Расположение атомов в упорядоченных сплавах показано на
рис.~\ref{fig:superlattices} (см. также п.~1.4.6). В CuAu атомы меди и
золота чередуются \emph{слоями}. Лучи, отражённые от основных и
промежуточных плоскостей (001), находятся в противофазе, но, рассеянные
разными атомами, гасят друг друга не полностью. На рентгенограммах
появляются линии, невозможные для ГЦК-решётки,~— \term{сверхструктурные
линии}. Поэтому упорядоченные растворы называют \term{сверхструктурами}.
Сверхструктурные линии тем ярче, чем сильнее различаются рассеивающие
способности атомов, т.\,е. их порядковые номера.

\begin{figure}[ht]
\centering
\begin{subfigure}[b]{0.3\textwidth}\centering
\begin{tikzpicture}[cubeproj, scale=1.9]
  \CubeEdges
  \foreach \p in {(0,1,0),(1,1,0),(0,1,1),(1,1,1)} \node[atomB=4.4mm] at \p {};
  \node[atom=3.6mm] at (0.5,1,0.5) {};
  \foreach \p in {(0,0.5,0.5),(1,0.5,0.5),(0.5,0.5,0),(0.5,0.5,1)} \node[atom=3.6mm] at \p {};
  \foreach \p in {(0,0,0),(1,0,0),(0,0,1),(1,0,1)} \node[atomB=4.4mm] at \p {};
  \node[atom=3.6mm] at (0.5,0,0.5) {};
\end{tikzpicture}
\caption{Cu$_3$Au}
\end{subfigure}\hfill
\begin{subfigure}[b]{0.3\textwidth}\centering
\begin{tikzpicture}[cubeproj, x={(1cm,0cm)}, y={(0.42cm,0.30cm)}, z={(0cm,0.935cm)}, scale=1.9]
  \CubeEdges
  \foreach \p in {(0,1,0),(1,1,0),(0,1,1),(1,1,1)} \node[atomB=4.4mm] at \p {};
  \node[atom=3.6mm] at (0.5,1,0.5) {};
  \foreach \p in {(0,0.5,0.5),(1,0.5,0.5)} \node[atom=3.6mm] at \p {};
  \foreach \p in {(0.5,0.5,0),(0.5,0.5,1)} \node[atomB=4.4mm] at \p {};
  \foreach \p in {(0,0,0),(1,0,0),(0,0,1),(1,0,1)} \node[atomB=4.4mm] at \p {};
  \node[atom=3.6mm] at (0.5,0,0.5) {};
\end{tikzpicture}
\caption{CuAu ($c/a=0{,}935$)}
\end{subfigure}\hfill
\begin{subfigure}[b]{0.3\textwidth}\centering
\begin{tikzpicture}[cubeproj, scale=1.9]
  \CubeEdges
  \foreach \p in {(0,1,0),(1,1,0),(0,1,1),(1,1,1)} \node[atom=4.0mm] at \p {};
  \node[atomI=4.0mm] at (0.5,0.5,0.5) {};
  \foreach \p in {(0,0,0),(1,0,0),(0,0,1),(1,0,1)} \node[atom=4.0mm] at \p {};
\end{tikzpicture}
\caption{CuZn ($\beta'$-латунь), FeAl}
\end{subfigure}
\caption{Упорядоченные твёрдые растворы: \textit{а}~— Cu$_3$Au (атомы Au,
оранжевые,~— в вершинах куба, Cu~— в центрах граней); \textit{б}~— CuAu:
слои Au и Cu чередуются вдоль оси $c$, решётка тетрагональна;
\textit{в}~— сверхструктура на основе ОЦК: атомы одного сорта в
вершинах, другого~— в центре куба}
\label{fig:superlattices}
\end{figure}

\begin{itemize}
\item Cu$_3$Au при упорядочении остаётся кубическим, а CuAu становится
\emph{тетрагональным} с укороченной осью $c$ ($c/a=0{,}935$): слои меди
«тоньше» слоёв золота, а период $a$ задают атомы золота. В системе
Cu--Au есть и тетрагональная сверхструктура Cu$_3$Au.
\item Ni$_3$Fe имеет ту же структуру, что Cu$_3$Au, но порядковые номера
Fe и Ni различаются лишь на 2, и сверхструктурные линии на
рентгенограммах не видны~— их обнаруживают только нейтронографически.
\item В CuPt атомы Cu и Pt заполняют чередующиеся плоскости (111).
\item Сверхструктуры на основе ОЦК~— сплавы Fe--Al (Fe$_3$Al и FeAl, см.
п.~1.4.6) и $\beta$-латунь CuZn (рис.~\ref{fig:superlattices},\,\textit{в});
на основе ГПУ~— MgCd$_3$, MgCd, Mg$_3$Cd.
\end{itemize}

\paragraph{Твёрдые растворы внедрения}
Атомы растворённого компонента занимают пустоты решётки. В ГЦК и ГПУ на
каждый атом приходятся одна октаэдрическая и две тетраэдрические
пустоты (п.~\ref{sssec:ionic}). В ОЦК-решётке пустоты тоже есть, но они
\emph{несимметричны}: в октаэдрическую пустоту вписывается не шар, а
сплюснутый эллипсоид с осями $0{,}16$ и $0{,}63$ диаметра атома
матрицы. На каждый атом ОЦК-решётки приходятся три октаэдрические и
шесть тетраэдрических пустот.

В междоузлия помещаются лишь атомы, диаметр которых не превышает
примерно половины диаметра атомов матрицы: H, C, N, O, B. Решётка и
металлические свойства основы сохраняются, период решётки растёт.
\vazhno{Неограниченной растворимости внедрения быть не может}: решётки
только из «внедрённых» атомов неметалла не существует. Обычно
растворимость мала, но иногда достигает атомной доли 0,25--0,30 (азот в
Ti, кислород в Zr), 0,5 (водород в Pd) и даже 0,75 (водород в Ce).

Тип занимаемой пустоты определяется тем же правилом, что и в ионных
кристаллах: внедрённый атом должен быть больше «своей» пустоты, но
меньше следующей по размеру.

\begin{primer}[Аустенит]
Аустенит~— раствор углерода в $\gamma$-Fe (ГЦК). Радиусы:
$r_\text{Fe}=0{,}126$~нм, $r_\text{C}=0{,}077$~нм, отношение $0{,}61$~—
в интервале $0{,}41$--$0{,}73$, поэтому углерод занимает
\emph{октаэдрические} пустоты (центр куба и середины рёбер ячейки). Это и
наблюдается; растворимость углерода в $\gamma$-Fe достигает 2,14~мас.\,\%,
т.\,е. около 9~ат.\,\%\isprmark. Радиус октаэдрической пустоты
$0{,}414\cdot0{,}126=0{,}052$~нм меньше радиуса углерода, поэтому решётка
расширяется (рис.~\ref{fig:a-vs-c},\,\textit{а}).
\end{primer}
\isprtext{В оригинале (с.~414)~— 8~ат.\,\%; предельная
растворимость углерода в аустените 2,14~мас.\,\% соответствует
9,2~ат.\,\%.}

По той же причине углерод и кислород в ОЦК-решётке железа, тантала,
ниобия и ванадия занимают октаэдрические пустоты, а водород в тантале и
ниобии~— тетраэдрические. Определить тип занятых пустот обычно удаётся
только нейтронографией: лёгкие атомы слабо рассеивают рентгеновские
лучи и электроны.
